\documentclass[11pt]{article}
\usepackage[a4paper,margin=1in]{geometry}
\usepackage{fontspec}
\usepackage[boldfont=false]{xeCJK}
\setCJKmainfont{FandolSong-Regular.otf}
\usepackage{amsmath}
\usepackage{amssymb}
\usepackage{array}
\usepackage{booktabs}
\usepackage{graphicx}
\usepackage{adjustbox}
\usepackage{enumitem}
\setlist[itemize]{topsep=2pt,partopsep=0pt,parsep=0pt,itemsep=2pt,leftmargin=1.4em}
\setlist[enumerate]{topsep=2pt,partopsep=0pt,parsep=0pt,itemsep=2pt,leftmargin=1.6em}
\usepackage{parskip}
\usepackage{fvextra}
\fvset{breaklines=true,breakanywhere=true,fontsize=\small}
\setlength{\parindent}{0pt}

\begin{document}

\begin{center}
  {\LARGE\bfseries Atoms, molecules and stoichiometry}\\[0.35em]
  {\large A Level Chemistry}
\end{center}
\vspace{0.6em}

\section*{Relative masses of atoms and molecules}

\textbf{Atoms} 原子 are far too light to weigh in grams, so we compare every mass to one standard. The standard is the \textbf{unified atomic mass unit} 统一原子质量单位 (symbol u), defined as exactly one twelfth of the mass of one carbon-12 atom.

Using this unit, we state masses as simple numbers:

\begin{itemize}
\item the \textbf{relative atomic mass} 相对原子质量 $A_r$ of an element is the average mass of its atoms compared with $\tfrac{1}{12}$ of a carbon-12 atom. It is an average over all the \textbf{isotopes} 同位素, weighted by how common each one is.

\item the \textbf{relative isotopic mass} 相对同位素质量 is the mass of one atom of a single isotope, compared with $\tfrac{1}{12}$ of a carbon-12 atom.

\item the \textbf{relative molecular mass} 相对分子质量 $M_r$ of a \textbf{molecule} 分子 is the sum of the relative atomic masses of all its atoms.

\item the \textbf{relative formula mass} 相对式量 is the same idea for a substance that is not made of molecules (such as an ionic compound). Add up the relative atomic masses shown in the formula.

\end{itemize}

\section*{The mole and the Avogadro constant}

Chemists count particles in groups called moles, just as we count eggs in dozens.

One \textbf{mole} 摩尔 (symbol mol) is the amount of substance that contains the same number of particles as there are atoms in exactly 12 g of carbon-12. That number is the \textbf{Avogadro constant} 阿伏伽德罗常量:

\[
N_A = 6.02 \times 10^{23}\ \text{mol}^{-1}
\]

So one mole of anything contains $6.02 \times 10^{23}$ particles. The particles may be atoms, molecules or \textbf{ions} 离子 — always say which.

The mass of one mole in grams equals the relative mass ($A_r$ or $M_r$). This is the \textbf{molar mass} 摩尔质量, with units $\text{g mol}^{-1}$. The key equation is:

\[
n = \frac{m}{M}
\]

where $n$ is the amount in moles, $m$ is the mass in grams, and $M$ is the molar mass.

\section*{Formulas}

A \textbf{compound} 化合物 is a substance made of two or more elements chemically joined.

\subsection*{Charges and formulas of ionic compounds}

In an \textbf{ionic compound} 离子化合物 the total positive charge balances the total negative charge, so the compound is neutral overall.

You can predict the charge of many ions from the element's position in the Periodic Table:

\begin{center}
\adjustbox{max width=\linewidth}{%
\begin{tabular}{lllllll}
\toprule
\textbf{Group} & \textbf{1} & \textbf{2} & \textbf{13} & \textbf{15} & \textbf{16} & \textbf{17} \\
\midrule
Usual ion charge & $+1$ & $+2$ & $+3$ & $-3$ & $-2$ & $-1$ \\
\bottomrule
\end{tabular}}
\end{center}

Hydrogen forms $\text{H}^+$, and the Group 18 noble gases do not normally form ions.

Some ions you must know by name and formula:

\begin{center}
\adjustbox{max width=\linewidth}{%
\begin{tabular}{ll}
\toprule
\textbf{Name} & \textbf{Formula} \\
\midrule
nitrate & $\text{NO}_3^{-}$ \\
carbonate & $\text{CO}_3^{2-}$ \\
sulfate & $\text{SO}_4^{2-}$ \\
hydroxide & $\text{OH}^{-}$ \\
ammonium & $\text{NH}_4^{+}$ \\
zinc & $\text{Zn}^{2+}$ \\
silver & $\text{Ag}^{+}$ \\
hydrogencarbonate & $\text{HCO}_3^{-}$ \\
phosphate & $\text{PO}_4^{3-}$ \\
\bottomrule
\end{tabular}}
\end{center}

For a metal that can have more than one charge, a Roman numeral shows the \textbf{oxidation number} 氧化数. For example, iron(II) is $\text{Fe}^{2+}$ and iron(III) is $\text{Fe}^{3+}$. To write a formula, balance the charges: iron(III) oxide is $\text{Fe}_2\text{O}_3$, because two $\text{Fe}^{3+}$ balance three $\text{O}^{2-}$.

\subsection*{Equations and state symbols}

A chemical equation must be \textbf{balanced} 配平 — the same number of each kind of atom on both sides. Add \textbf{state symbols} 状态符号 to show the state of each species: (s) solid, (l) liquid, (g) gas, and (aq) \textbf{aqueous} 水溶液 (dissolved in water).

An \textbf{ionic equation} 离子方程式 shows only the ions and molecules that actually change. The ions that do not change are \textbf{spectator ions} 旁观离子, and you leave them out. For example, the reaction that forms silver chloride is:

\[
\text{Ag}^{+}(\text{aq}) + \text{Cl}^{-}(\text{aq}) \rightarrow \text{AgCl}(\text{s})
\]

\subsection*{Empirical and molecular formulas}

The \textbf{empirical formula} 实验式 is the simplest whole-number ratio of the atoms of each element in a compound. The \textbf{molecular formula} 分子式 shows the actual number of atoms of each element in one molecule.

For example, ethane has empirical formula $\text{CH}_3$ but molecular formula $\text{C}_2\text{H}_6$.

\subsection*{Hydrated and anhydrous solids}

Some solids hold water inside their crystals. This water is the \textbf{water of crystallisation} 结晶水. A solid that contains it is \textbf{hydrated} 水合的; the same solid with the water removed is \textbf{anhydrous} 无水的.

For example, hydrated copper(II) sulfate is $\text{CuSO}_4 \cdot 5\text{H}_2\text{O}$. Heating it drives off the water to leave anhydrous $\text{CuSO}_4$.

\subsection*{Calculating empirical and molecular formulas}

To find the empirical formula from masses (or percentages by mass):

\begin{enumerate}
\item divide each element's mass by its $A_r$ to get the moles.

\item divide all the mole values by the smallest one.

\item round to the nearest whole numbers — that ratio is the empirical formula.

\end{enumerate}

To get the molecular formula, you also need $M_r$. Find how many times the empirical formula mass fits into $M_r$, then multiply the formula by that number.

\section*{Reacting masses and volumes}

\subsection*{Reacting masses and percentage yield}

The numbers in front of each species in a balanced equation give the mole ratio — this is the \textbf{stoichiometry} 化学计量. To find a reacting mass: change the known mass to moles, use the mole ratio to find the moles you want, then change back to mass.

In real reactions you usually get less product than the maximum. The \textbf{percentage yield} 产率 compares the amount you actually made with the most you could make:

\[
\text{percentage yield} = \frac{\text{actual amount of product}}{\text{maximum possible amount}} \times 100\%
\]

\subsection*{Limiting and excess reagent}

When two reactants are mixed, one usually runs out first. The \textbf{limiting reagent} 限量试剂 is the one that runs out — it decides how much product forms. The other is the \textbf{excess reagent} 过量试剂, because there is more than enough of it. Always base the calculation on the limiting reagent.

To find it: work out the moles of each reactant, divide each by its number in the equation, and the smallest result is the limiting reagent.

\subsection*{Volumes of gases}

At the same temperature and pressure, equal volumes of any gases contain equal numbers of molecules. At room temperature and pressure (r.t.p.), one mole of any gas takes up $24.0\ \text{dm}^3$, so:

\[
n = \frac{V}{24.0}\qquad (V \text{ in } \text{dm}^3 \text{ at r.t.p.})
\]

This is used when burning \textbf{hydrocarbons} 碳氢化合物 (compounds of only carbon and hydrogen), where you compare gas volumes.

\subsection*{Volumes and concentrations of solutions}

The \textbf{concentration} 浓度 of a solution is the amount of \textbf{solute} 溶质 in each cubic decimetre of \textbf{solution} 溶液, measured in $\text{mol dm}^{-3}$:

\[
n = c \times V
\]

where $c$ is the concentration and $V$ is the volume in $\text{dm}^3$. Remember that $1000\ \text{cm}^3 = 1\ \text{dm}^3$.

This is the basis of a \textbf{titration} 滴定, where you find an unknown concentration by reacting it with a solution whose concentration you already know.

\subsection*{Significant figures}

Give your answer to a sensible number of \textbf{significant figures} 有效数字 — usually match the data in the question. Do not write more digits than the data supports, and do not round so early that you lose accuracy.


\end{document}
